\chapter{Acetylcholine Receptor Antibodies}

\section*{What Is This Test?}
The acetylcholine receptor (AChR) antibody test measures autoantibodies directed against the nicotinic acetylcholine receptor at the neuromuscular junction. These antibodies are the diagnostic marker of myasthenia gravis (MG). MG is an autoimmune disorder of the neuromuscular junction. It causes fluctuating weakness of skeletal muscles. This is classically ocular (ptosis, diplopia) and bulbar (dysarthria, dysphagia). It causes fatigability that worsens with use and improves with rest. In most patients with generalized MG, autoantibodies to the acetylcholine receptor are detectable in serum. The antibodies are divided into three subtypes. Binding antibodies are directed against the receptor. They cause receptor blockade and complement-mediated destruction. Blocking antibodies block the binding of acetylcholine. Modulating antibodies accelerate receptor internalization. The most commonly ordered test is the AChR binding antibody assay. A positive result confirms the diagnosis of seropositive MG. A negative result in a patient with clinically suspected MG does not exclude the diagnosis. This is seronegative MG. In that case, testing for MuSK, LRP4, and agrin antibodies is pursued. Electrophysiologic studies are also pursued. Approximately 10--15\% of patients with generalized MG are seronegative for AChR antibodies. Up to 50\% of patients with purely ocular MG are seronegative \cite{gilhus2016myasthenia}.

\section*{Alternative Names}
\begin{itemize}
  \item Acetylcholine Receptor Antibody Test
  \item Anti-AChR Antibodies
  \item AChR Ab
  \item Acetylcholine Receptor Binding Antibodies
  \item AChR Binding, Blocking, and Modulating Antibodies
  \item Myasthenia Gravis Antibody Panel
\end{itemize}
\section*{Principle}
The AChR binding antibody assay uses a radioimmunoprecipitation technique. A purified acetylcholine receptor is labeled with iodine-125 alpha-bungarotoxin. The receptor typically comes from human muscle or a stable cell line expressing the receptor. Alpha-bungarotoxin is a snake toxin. It binds specifically and irreversibly to the acetylcholine receptor. It binds at a site distinct from the acetylcholine binding site. Patient serum is incubated with the radiolabeled receptor. If AChR antibodies are present, they bind to the receptor. Anti-human immunoglobulin is then added to precipitate the immune complexes. The radioactivity of the precipitate is counted. The amount of precipitated radioactivity is proportional to the concentration of AChR binding antibodies. Results are reported in nmol/L. Blocking antibodies are measured by assessing whether patient serum inhibits the binding of a labeled monoclonal antibody to the toxin-binding site. Modulating antibodies are measured by the ability of patient serum to accelerate receptor degradation in a cultured muscle cell line. Binding antibodies are the most sensitive. They are the first-line test. Blocking and modulating assays are added when the binding assay is negative but MG is strongly suspected. The presence of AChR antibodies is specific. A positive result strongly supports MG. In a minority of patients, they may also be found in other conditions. These include Lambert-Eaton myasthenic syndrome, autoimmune hepatitis, and thymoma without MG.

\section*{History}
Myasthenia gravis was first described in detail by Thomas Willis in the 17th century. The modern clinical entity was defined by Jolly in 1895. Jolly also coined the term ``myasthenia gravis pseudoparalytica.'' The autoimmune basis was established in 1973 by Patrick and Lindstrom. They induced experimental autoimmune MG in rabbits by immunizing them with purified acetylcholine receptor. This demonstrated that the receptor is the target antigen \cite{patrick1973autoimmune}. In 1976, Lindstrom and colleagues detected antibodies against the acetylcholine receptor in the serum of patients with MG. This established the diagnostic test \cite{lindstrom1976antibody}. The radioimmunoprecipitation assay using alpha-bungarotoxin became the standard clinical assay. Subsequent work identified the MuSK (muscle-specific kinase) antibodies in 2001. Hoch and Vincent described them. This explained a major subset of seronegative MG \cite{hoch2001autoantibodies}. The role of the thymus in MG pathogenesis was recognized. Thymoma occurs in 10--15\% of patients. Thymic hyperplasia occurs in most others. So chest imaging and thymectomy became standard components of management. The MGTX trial published in 2016 validated this \cite{wolfe2016randomized}.

\section*{Why This Test Is Ordered}
\begin{itemize}
  \item \textbf{Evaluation of suspected myasthenia gravis:} In patients with fluctuating ptosis, diplopia, bulbar weakness (dysarthria, dysphagia, choking), or fatigable proximal limb weakness
  \item \textbf{Diagnosis of ocular myasthenia:} Isolated ocular symptoms; note that up to 50\% of ocular MG is seronegative for AChR antibodies
  \item \textbf{Distinguishing MG from other neuromuscular junction and motor disorders:} Lambert-Eaton myasthenic syndrome, congenital myasthenic syndromes, amyotrophic lateral sclerosis, and botulism
  \item \textbf{Evaluation before and after thymectomy:} To document the antibody status and, after surgery, to monitor improvement
  \item \textbf{Monitoring response to immunotherapy:} Antibody titers often fall with successful treatment, although titers correlate imperfectly with clinical status
  \item \textbf{Screening in patients with thymoma:} MG occurs in approximately one-third of patients with thymoma
\end{itemize}

\section*{Primary vs Secondary Test}
The AChR antibody test is the primary laboratory test for the diagnosis of myasthenia gravis. It is a secondary step in the diagnostic pathway. The diagnosis begins with a careful history and neurologic examination. It is confirmed by serology plus electrodiagnostic studies. These studies include repetitive nerve stimulation and single-fiber EMG. When appropriate, the edrophonium test is used. Because a negative result does not exclude MG, the test is interpreted with the full clinical picture.

\section*{How This Test Is Performed}
A venous blood sample is collected in a red-top (serum separator, SST) tube. It is typically 3--5 mL. No special preparation is required. The serum is sent to a reference laboratory for the radioimmunoprecipitation assay for AChR binding antibodies. If the binding assay is negative and clinical suspicion remains high, blocking and modulating antibody assays are added. MuSK and LRP4 antibodies are added as well. Results are typically available within 3--7 days. The quantitative result is reported in nmol/L. Values above the laboratory's reference cutoff are considered positive (commonly $>$0.4--0.5 nmol/L, though cutoffs vary by laboratory).

\section*{What Preparation Is Needed}
No fasting is required. The patient should inform the healthcare provider of all medications. This especially includes cholinesterase inhibitors (pyridostigmine, which do not interfere with the antibody assay but should be continued until instructed). It also includes corticosteroids and immunosuppressants. These can lower antibody titers and cause false negatives if started before testing. The patient should report any history of thymoma or other autoimmune disease. The test should ideally be performed before immunosuppressive therapy is initiated.

\section*{Contraindications and Precautions}
\begin{itemize}
  \item There are no absolute contraindications to the blood draw. Important interpretive cautions:
  \item (1) a negative AChR antibody result does not exclude MG --- up to 50\% of ocular MG and 10--15\% of generalized MG are seronegative; MuSK and LRP4 antibodies and electrodiagnostic testing are required;
  \item (2) AChR antibodies can be positive in Lambert-Eaton myasthenic syndrome, autoimmune hepatitis, and some patients with thymoma without clinical MG, so results must be interpreted clinically;
  \item (3) antibody titers correlate imperfectly with clinical severity and should not be used alone to adjust therapy;
  \item (4) immunosuppressive therapy, including steroids and rituximab, can lower titers and produce false negatives;
  \item (5) the radioimmunoprecipitation assay uses radioactivity, but no radioactive material is given to the patient;
  \item (6) seroconversion from seronegative to seropositive may occur over time, so repeat testing is reasonable in initially negative patients whose clinical picture evolves;
  \item (7) the presence of anti-AChR modulating antibodies alone is rare, and the panel is most informative when all three subtypes are tested.
\end{itemize}
\section*{Risks}
Risks are those of routine venipuncture. They include mild pain, bruising, and small hematoma at the collection site. They also include lightheadedness or vasovagal syncope. Extremely rare risks are infection or nerve injury.

\section*{What Happens After the Test}
If the AChR antibody test is positive, the diagnosis of seropositive MG is confirmed. The patient is referred to a neurologist. Chest imaging (computed tomography of the chest) is performed to evaluate for thymoma and thymic hyperplasia. Treatment is initiated with a cholinesterase inhibitor (pyridostigmine) for symptomatic control. Corticosteroids and other immunosuppressants are added for moderate to severe disease. These include azathioprine, mycophenolate, and rituximab. Thymectomy is recommended for patients with thymoma. It is also recommended for many patients with generalized AChR antibody-positive MG without thymoma. The MGTX trial showed benefit in patients under 65 years with generalized nonthymomatous MG. If the test is negative, the workup proceeds with MuSK and LRP4 antibody testing. It also proceeds with repetitive nerve stimulation, single-fiber EMG, and a trial of cholinesterase inhibitors. Alternative diagnoses are considered. These include Lambert-Eaton myasthenic syndrome, congenital myasthenic syndromes, botulism, and motor neuron disease \cite{sanders2016international}. MG can be life-threatening (myasthenic crisis with respiratory failure). So a positive diagnosis prompts education about warning signs. It also prompts education about medications that can worsen MG (certain antibiotics, magnesium, beta-blockers).

\section*{Understanding Results}

\subsection*{Positive (AChR Antibodies Detected)}
\begin{itemize}
  \item \textbf{Myasthenia gravis (seropositive):} Confirms the diagnosis in the appropriate clinical context; antibody titers are often higher in generalized than in ocular MG
  \item \textbf{Thymoma:} Approximately 10--15\% of MG patients have a thymoma; conversely, about one-third of thymoma patients have MG and may have AChR antibodies even without clinical disease
  \item \textbf{Other autoimmune conditions:} Rarely positive in Lambert-Eaton myasthenic syndrome, autoimmune hepatitis, and other autoimmune diseases without MG
  \item \textbf{Persistent after treatment:} Titers often decline with immunosuppression and after thymectomy, but clinical correlation is imperfect and the value is used mainly to confirm the diagnosis
\end{itemize}

\subsection*{Negative}
\begin{itemize}
  \item \textbf{Seronegative MG:} Generalized MG (10--15\%) or ocular MG (up to 50\%); testing for MuSK, LRP4, and agrin antibodies is pursued
  \item \textbf{Acquired neuromyotonia and other NMJ disorders:} Different antibody targets (CASPR2, etc.)
  \item \textbf{Alternative neuromuscular disorders:} Lambert-Eaton myasthenic syndrome (antibodies to voltage-gated calcium channels), congenital myasthenic syndromes (genetic), botulism, or motor neuron disease
  \item \textbf{False negative:} Testing after immunosuppressive therapy has lowered titers, or early disease with low titers
\end{itemize}

\section*{Normal Results}
\begin{itemize}
  \item \textbf{AChR binding antibodies:} Negative (below the laboratory cutoff, typically $<$0.4--0.5 nmol/L; some laboratories report $<$0.02--0.5 nmol/L depending on assay)
  \item \textbf{AChR blocking and modulating antibodies:} Negative (low or absent)
  \item \textbf{Positive threshold:} Values above the laboratory's reference cutoff are reported as positive; moderately elevated values (e.g., 0.5--10 nmol/L) are common in generalized MG, and very high values are seen in some patients with severe disease or thymoma
\end{itemize}

\section*{Follow-up Tests}
\begin{itemize}
  \item \textbf{MuSK (muscle-specific kinase) antibodies:} In seronegative MG, especially with prominent bulbar weakness
  \item \textbf{LRP4 and agrin antibodies:} In double-seronegative MG
  \item \textbf{Repetitive nerve stimulation:} Decremental response at low frequencies supports MG
  \item \textbf{Single-fiber electromyography:} The most sensitive electrodiagnostic test for MG
  \item \textbf{Edrophonium (Tensilon) test:} Historical and selective adjunct for diagnosis
  \item \textbf{Chest computed tomography or MRI:} To detect thymoma and thymic hyperplasia
  \item \textbf{Anti-striational antibodies:} Associated with thymoma and late-onset MG
  \item \textbf{Voltage-gated calcium channel antibodies:} When Lambert-Eaton myasthenic syndrome is suspected
  \item \textbf{Pulmonary function testing:} For bulbar and respiratory weakness in severe disease
\end{itemize}

\section*{References}
\bibliographystyle{unsrtnat}
\bibliography{references}

\clearpage
\section*{Regulatory Status and Authorization Date}
This chapter describes a test category, not a specific commercial product. FDA, CDSCO, EU IVDR, and MHRA approval or registration dates therefore cannot be assigned without the assay manufacturer, product name, instrument, and intended use. For laboratory-developed tests, an individual agency approval date may not exist. See the Preface for the interpretation of regulatory status.

\clearpage
